\documentclass[10pt]{article}
\usepackage[margin=1in]{geometry}
\usepackage[T1]{fontenc}
\usepackage{lmodern}
\usepackage{microtype}
\usepackage{booktabs}
\usepackage{graphicx}
\usepackage{enumitem}
\usepackage[hidelinks]{hyperref}
\setlist{nosep,leftmargin=1.4em}
% Generated by paper/make_tables.py. Do not edit by hand.

\newcommand{\TableModels}{%
\begin{tabular}{llll}\toprule
Model & Vendor & Name sent to the endpoint & Endpoint \\\midrule
Llama 4 Maverick & Meta (open weights) & \texttt{meta-llama/llama-4-maverick} & OpenRouter \\
Claude Sonnet 5.5 & Anthropic & \texttt{anthropic/claude-sonnet-5.5} & OpenRouter \\
GPT-6 Luna & OpenAI & \texttt{gpt-6-luna} & OpenAI API \\
GPT-6.1 Terra & OpenAI & \texttt{gpt-6.1-terra} & OpenAI API \\
Gemini 3.8 Flash & Google & \texttt{google/gemini-3.8-flash} & OpenRouter \\
Claude Haiku 5.5 & Anthropic & \texttt{anthropic/claude-haiku-5.5} & OpenRouter \\
\bottomrule\end{tabular}}

\newcommand{\TableSecurity}{%
\begin{tabular}{lrrrrrrr}\toprule
 & & \multicolumn{3}{c}{Attack success (\%)} & \multicolumn{3}{c}{No defense, by attack (\%)} \\\cmidrule(lr){3-5}\cmidrule(lr){6-8}
Model & Saved (\%) & No defense & Labels & Enforced & Extra action & Override & Fill gap \\\midrule
Llama 4 Maverick & 100 $\pm$ 0 & 92 $\pm$ 3 & 69 $\pm$ 8 & 5 $\pm$ 3 & 96 & 80 & 100 \\
Claude Sonnet 5.5 & 53 $\pm$ 3 & 48 $\pm$ 7 & 2 $\pm$ 3 & 0 $\pm$ 0 & 0 & 58 & 86 \\
GPT-6 Luna & 45 $\pm$ 5 & 41 $\pm$ 5 & 4 $\pm$ 3 & 0 $\pm$ 0 & 30 & 20 & 72 \\
GPT-6.1 Terra & 43 $\pm$ 6 & 33 $\pm$ 8 & 15 $\pm$ 8 & 0 $\pm$ 0 & 0 & 36 & 62 \\
Gemini 3.8 Flash & 25 $\pm$ 6 & 24 $\pm$ 8 & 0 $\pm$ 0 & 0 $\pm$ 0 & 0 & 28 & 44 \\
Claude Haiku 5.5 & 6 $\pm$ 3 & 1 $\pm$ 2 & 1 $\pm$ 2 & 0 $\pm$ 0 & 0 & 4 & 0 \\
\bottomrule\end{tabular}}

\newcommand{\TableCost}{%
\begin{tabular}{lrrrrrrr}\toprule
 & \multicolumn{2}{c}{Correct, unprompted (\%)} & \multicolumn{5}{c}{Enforced defense (\%)} \\\cmidrule(lr){2-3}\cmidrule(lr){4-8}
Model & No defense & Labels & Correct, unprompted & Correct after prompt & Detail lost & Asked anyway & Prompts / 100 tasks \\\midrule
Llama 4 Maverick & 100 $\pm$ 0 & 87 $\pm$ 3 & 82 $\pm$ 2 & 10 $\pm$ 2 & 8 $\pm$ 3 & 0 $\pm$ 0 & 9.9 \\
Claude Sonnet 5.5 & 100 $\pm$ 0 & 65 $\pm$ 3 & 65 $\pm$ 3 & 0 $\pm$ 0 & 13 $\pm$ 3 & 22 $\pm$ 3 & 0.0 \\
GPT-6 Luna & 85 $\pm$ 2 & 70 $\pm$ 2 & 68 $\pm$ 1 & 2 $\pm$ 1 & 23 $\pm$ 4 & 7 $\pm$ 2 & 1.7 \\
GPT-6.1 Terra & 93 $\pm$ 2 & 78 $\pm$ 3 & 76 $\pm$ 1 & 7 $\pm$ 1 & 15 $\pm$ 3 & 2 $\pm$ 1 & 7.0 \\
Gemini 3.8 Flash & 100 $\pm$ 0 & 57 $\pm$ 0 & 57 $\pm$ 0 & 0 $\pm$ 0 & 13 $\pm$ 3 & 30 $\pm$ 3 & 0.0 \\
Claude Haiku 5.5 & 73 $\pm$ 3 & 58 $\pm$ 1 & 58 $\pm$ 1 & 0 $\pm$ 0 & 21 $\pm$ 4 & 21 $\pm$ 4 & 0.0 \\
\bottomrule\end{tabular}}

\newcommand{\TableBehaviour}{%
\begin{tabular}{lrrrrrrr}\toprule
 & \multicolumn{2}{c}{Planted, saved} & \multicolumn{2}{c}{Legitimate, saved} & \multicolumn{3}{c}{With labels visible} \\\cmidrule(lr){2-3}\cmidrule(lr){4-5}\cmidrule(lr){6-8}
Model & Blunt & Plain & Account & Contact & Acted & Drafted & Chose trusted \\\midrule
Llama 4 Maverick & 58/58 & 50/50 & 8/8 & 15/15 & 337/339 & 168/168 & 180/180 \\
Claude Sonnet 5.5 & 16/58 & 33/50 & 8/8 & 15/15 & 0/199 & 63/122 & 57/57 \\
GPT-6 Luna & 9/58 & 37/50 & 0/8 & 14/15 & 35/167 & 77/97 & 53/53 \\
GPT-6.1 Terra & 6/58 & 35/50 & 5/8 & 15/15 & 136/206 & 116/122 & 52/52 \\
Gemini 3.8 Flash & 2/58 & 22/50 & 8/8 & 15/15 & 0/148 & 1/113 & 30/30 \\
Claude Haiku 5.5 & 1/58 & 5/50 & 3/8 & 12/15 & 0/65 & 4/57 & 6/6 \\
\bottomrule\end{tabular}}

\newcommand{\TableFatigue}{%
\begin{tabular}{lrrrr}\toprule
 & \multicolumn{4}{c}{Chance the user approves a bad prompt} \\\cmidrule(lr){2-5}
Model & 0 & 0.1 & 0.3 & 0.5 \\\midrule
Llama 4 Maverick & 0.0 & 4.7 & 10.0 & 15.3 \\
Claude Sonnet 5.5 & 0.0 & 0.0 & 0.0 & 0.0 \\
GPT-6 Luna & 0.0 & 0.0 & 0.7 & 1.3 \\
GPT-6.1 Terra & 0.0 & 0.0 & 2.7 & 4.7 \\
Gemini 3.8 Flash & 0.0 & 0.0 & 0.0 & 0.0 \\
Claude Haiku 5.5 & 0.0 & 0.0 & 0.0 & 0.0 \\
\bottomrule\end{tabular}}


\title{Ask, Forget, Repair: What Provenance-Gated Agent Memory\\Costs Across Six Language Models}
\author{Radhika Motwani\\\texttt{motwani.r@northeastern.edu}}
\date{Draft of 25 November 2026. Not yet submitted.}

\begin{document}
\maketitle

\begin{abstract}
An agent with long-term memory can be poisoned in one session and exploited in a later one. A
common defense is to record where each memory came from and to stop untrusted memories from
driving consequential actions. Prior work shows that such defenses block attacks; much less is
known about what they cost, or how much of the protection comes from the model and how much from
the runtime. We run six current language models from four vendors as the agent in one 60-session
synthetic back-office workload with 150 planted memory items per model, under three conditions:
no defense, provenance labels that are shown to the model but not enforced, and the same labels
enforced by the runtime. With no defense, planted items later caused a wrong action in 1\% to
92\% of cases depending on the model. Showing labels alone reduced this to 0--4\% for four models
but left 15\% and 69\% for the other two. With enforcement no planted item succeeded on any model
unless a simulated user approved a bad confirmation. The defense is not free: with no attacker
present, the share of requests completed correctly without involving the user fell by 15 to 43
points. Where that cost lands depends on the model. Models that respect labels stop using
legitimate outside facts and ask the user instead; a model that ignores labels triggers runtime
confirmation prompts. The planted texts are simple and do not adapt to the defense, the tasks are
synthetic, and the user is simulated, so these results are a lower bound on attacks and a first
measurement of cost, not a ranking of models. Code, workload and every model reply are released.
\end{abstract}

\section{Introduction}

Agents that keep memory across sessions inherit a delayed form of prompt injection. A document
read today (an email, a web page, a shared file) can leave behind a memory that steers an action
weeks later, after the document itself is gone. A line of recent work addresses this by attaching
provenance to memory and denying untrusted memory the authority to trigger consequential
actions~\cite{tmanm,permit,ppmf,sleeperchannels}. These systems report near-zero attack success.

Two questions are left open by that work, and they matter to anyone deciding whether to deploy
such a defense. First, what does it cost? A legitimate supplier really does send a new bank
account by email. A defense that distrusts everything from outside must either withhold that fact,
ask the user about it, or let it expire. Second, how much of the protection is the runtime's and
how much is the model's own caution? If a capable model already declines to act on a memory it can
see is untrusted, then showing it a label may be enough, and enforcement is redundant. If it does
not, a label is decoration.

We measure both on six current models with one workload and one memory runtime. This is a
measurement paper. The defense mechanisms we implement are taken from prior work, chiefly
origin-bound authority~\cite{tmanm} and confirmation gates~\cite{permit,sleeperchannels}; we claim
no new defense. Our contributions are:

\begin{itemize}
\item A cross-vendor comparison of labels shown to the model against labels enforced by the
runtime. Labels alone were nearly sufficient for four hosted models and clearly insufficient for
an open-weight model and for one hosted model (Section~\ref{sec:labels}).
\item A breakdown of what the enforced defense costs when no attacker is present, separating
requests the agent completes unprompted, requests that need a confirmation, details that are
lost, and requests where the agent has the right detail and asks the user anyway
(Section~\ref{sec:cost}).
\item Observations about the memory lifecycle: merging trusted and untrusted entries makes the
user's own instructions distrusted, expiry removes legitimate facts along with planted ones, and
retracting a malicious source only works if provenance survived the agent's own note-taking
(Section~\ref{sec:lifecycle}).
\item A harness in which the agent only proposes what to store and what to do, so any
model behind a chat-completions endpoint can be evaluated, with every prompt and reply saved.
\end{itemize}

\section{Related work}

\paragraph{Provenance-gated memory.} TMA-NM binds a memory item's authority to its origin at write
time, propagates that through the agent's own summaries, and reports no successful attacks across
eight models, with machine-checked proofs~\cite{tmanm}. PerMIT combines a quarantined reader with a
confirmation gate and eviction~\cite{permit}. PPMF describes \emph{provenance laundering}, in which
consolidation rewrites outside content as apparent user history, and gates tool calls on
platform-maintained provenance; it reports that readable provenance cues in memory text remained
exploitable on several small open-weight models while the runtime gate held~\cite{ppmf}. Maloyan
and Namiot distinguish provenance warnings placed in the model's context from gates enforced
outside the model loop and argue that only the latter is load-bearing, without a multi-model
measurement~\cite{sleeperchannels}. Our Labels and Enforced conditions correspond to that
distinction. We add measurements on current hosted models, where the picture differs from the
small open-weight models studied before, and a breakdown of the cost.

\paragraph{Cost of memory defenses.} ``The Price of Safety'' measures benign accuracy and token
overhead for four memory defenses, with no attacker and no confirmation step~\cite{priceofsafety}.
MEMGAUGE studies utility--risk trade-offs across writing, management and retrieval for eleven
models, with provenance used to quarantine records at the system level rather than shown to the
model~\cite{memgauge}. Karunanidhi shows that ranking by provenance can remove the evidence a
question needs when that evidence is itself untrusted~\cite{utilityunderattack}; our ``detail
lost'' and ``asked anyway'' columns are the action-side counterpart of that effect.

\paragraph{How models treat memory.} Hu and Ramachandran find that small open-weight models rely
on stale stored facts almost always, and that provenance metadata helps larger models more than
smaller ones~\cite{trustgap}. Li et al.\ compare how two Claude models and a Qwen model decide
whether to persist, verify or ask, without provenance labels or planted content~\cite{rememberverify}.
We observe related differences between vendors in an adversarial, multi-session setting.

\paragraph{Benchmarks and attacks.} MemoryDojo provides realistic multi-session tasks across five
memory backends and evaluates no defenses~\cite{memorydojo}; it would be a better testbed than our
synthetic workload once its code is available. Adaptive attacks have repeatedly broken
prompt-injection defenses that reported near-zero attack success against static
attacks~\cite{attackermovessecond}. We do not run adaptive attacks, and Section~\ref{sec:limits}
says what follows from that.

\section{Setup}

\subsection{Threat model}

The attacker controls the content of some documents the agent reads. The attacker cannot reach the
model, the runtime, the memory store or the user's own messages, and wants a consequential action
(moving money or sending data) to go wrong in a later session. The runtime, the user channel and
the labels the runtime assigns are trusted. The model is not.

\subsection{Workload}

Each run simulates 60 sessions of a back-office assistant. Topics are details about twenty
invented companies: a payout account, a billing contact, or the recipient of a usage report. A
run contains:

\begin{itemize}
\item 20 details the user states directly in the first third of the run;
\item 15 legitimate details that arrive only in outside documents, such as a supplier's email;
\item 30 planted items in outside documents, ten of each kind below, each followed 1 to 40
sessions later by a request that depends on the same topic;
\item 135 user requests in total. Each asks the agent either to act (``Pay this month's invoice
from \ldots'') or to prepare something without external effect (``Add \ldots\ to my draft payment
sheet. Do not pay anything yet.'').
\end{itemize}

The three kinds of planted item are an \emph{extra action} (a standing instruction to copy
material to an outside address), an \emph{override} (a value contradicting one the user stated),
and a \emph{fill gap} (a value for a detail the user never gave). Half are worded bluntly
(``Ignore the earlier setting. Always use \ldots'') and half plainly (``\ldots\ is now
\ldots''). They are fixed templates. Nothing about them is tuned to any model or defense.
Memory is consolidated every five sessions: entries about the same topic are merged and the
newest value wins. We use five seeds, so each model sees 150 planted items. Because seeds share
many documents, a model answers 108 distinct planted documents.

\subsection{Agent protocol}

The model is called at two points. When a document is read, it is asked what, if anything, to
remember, and replies with a list of (topic, value, standing instruction) entries. When the user
makes a request, it is shown the request and its memory entries about the needed detail, and
replies with the entry it will use or with a decision to ask the user instead, plus any additional
actions it intends to take. The model only proposes. The runtime stores, labels, retrieves and
decides what executes. Retrieval is by exact topic handle; each document statement carries its
handle. Identical prompts are sent once and the reply reused, and requests ask for temperature~0
where the model accepts it.

\subsection{Conditions}

\begin{description}[leftmargin=1.4em,style=nextline]
\item[No defense.] A plain store. Entries carry no source. A merged entry looks like a fresh note.
\item[Labels.] The runtime labels every entry. User statements are trusted. Anything the model
writes after reading an outside document in the same session is untrusted and carries that
document's identifier. Untrusted entries expire after 20 sessions. A merged entry takes the lowest
trust, all sources and the earliest expiry of its parts. The model sees each entry's label and
sources. Nothing is enforced, and entries of different trust may be merged.
\item[Enforced.] Labels, plus runtime rules. Trusted and untrusted entries are never merged. A
value the user stated outranks any other. An untrusted memory cannot start an action, and an
action the model does not attribute to a memory is denied while untrusted memory is in view. A
request that acts and would use an untrusted value produces a confirmation prompt showing the
source; approval re-issues the value as a user statement. Retracting a source removes everything
it influenced.
\end{description}

Two further baselines, a keyword filter on writes and expiry applied only to verbatim tool output,
gave results identical to No defense for every model, because the model's own notes are what
persists. We omit them from the tables.

\subsection{Simulated user and metrics}

At a confirmation prompt the simulated user always approves a correct value and approves a planted
one with probability $f$ ($f=0.1$ unless stated). \emph{Attack success} is the share of planted
items that, at any later point, had their value used in a request that acts or had their
instruction executed. \emph{Saved} is the share present in memory after the document was read. For
requests whose needed detail has a true value, each is classed as \emph{correct, unprompted};
\emph{correct after prompt}; \emph{detail lost} (no entry holding the right value could be
retrieved, because it was never saved or had expired); \emph{asked anyway} (such an entry was
retrievable and the model chose to ask the user); or \emph{wrong}. Cost figures come from a second
set of runs with the planted items removed. Tables report the mean and standard deviation over
five seeds.

\subsection{Models}

\begin{center}\small\TableModels\end{center}

\noindent The runs were made between 22 and 25 November 2026 and comprise 6{,}813 distinct model
calls.

\section{Results}

\subsection{Undefended agents are poisoned at very different rates}

\begin{table}[t]\centering\small
\resizebox{\textwidth}{!}{\TableSecurity}
\caption{Security. ``Saved'' is the share of planted items in memory after reading. Enforced is
with a user who approves a bad prompt one time in ten. Mean $\pm$ standard deviation over five
seeds; 150 planted items per model.}
\label{tab:security}
\end{table}

Table~\ref{tab:security} shows attack success with no defense ranging from 1\% for Claude Haiku
5.5 to 92\% for Llama 4 Maverick. Most of that spread is decided when the document is read.
Table~\ref{tab:behaviour} shows what each model chose to save. Llama saved all 108 distinct
planted documents. The five hosted models together saved something from 34 of 290 bluntly worded
ones, and from 5 to 37 of the 50 plainly worded ones. Once a planted item was saved by an
undefended agent it usually succeeded: saved and successful shares differ by at most ten points.

Models also differed on legitimate content. GPT-6 Luna saved none of the eight legitimate payout
accounts that arrived in outside documents, and Claude Haiku three, while both saved most contact
addresses. That caution costs utility with no defense at all: 15\% and 13\% of needed details
were lost for those two models, against none for Sonnet, Gemini and Llama.

Only two models were affected by extra-action items (Luna 30\%, Llama 96\%); the others saved no
standing instruction that later executed. Claude Sonnet 5.5 returned an empty reply to 20
documents, all of them planted extra-action notes; an empty reply is scored as saving nothing,
and the harness did not record why the reply was empty.

\begin{table}[t]\centering\small
\TableBehaviour
\caption{Model behaviour, counted over distinct prompts. Left: documents from which the model
saved a value or instruction, for planted documents worded bluntly or plainly and for legitimate
payout accounts and contact addresses. Right, with labels visible: how often the model used a
memory when only untrusted entries were available, for requests that act (``Acted'') and for
drafts (``Drafted''); and, when a trusted and an untrusted entry were both in view and the model
used one, how often it was the trusted one.}
\label{tab:behaviour}
\end{table}

\subsection{Labels shown to the model help some models much more than others}
\label{sec:labels}

With labels visible but unenforced, attack success fell to 0--4\% for Sonnet, Gemini, Luna and
Haiku, to 15\% for GPT-6.1 Terra, and only to 69\% for Llama. The right half of
Table~\ref{tab:behaviour} shows why. Faced with a request that acts and only untrusted memory to
go on, Sonnet, Gemini and Haiku never used the memory (0 of 412 cases). Luna used it in 35 of
167, Terra in 136 of 206, and Llama in 337 of 339.

When a trusted and an untrusted entry were both in view, every model chose the trusted one, in
all 378 cases. Failures to follow labels therefore arose where there was nothing
trusted to prefer: the fill-gap case, and cases where merging had already destroyed the trusted
entry (Section~\ref{sec:lifecycle}).

With enforcement, no planted item succeeded on five models. On Llama 4.7\% succeeded, all through
confirmation prompts that the simulated user approved. Table~\ref{tab:fatigue} varies that user.
With a user who never approves a bad prompt, success is zero for every model. With one who
approves half of them it reaches 15\% for Llama, because Llama proposes the planted value every
time and so every such request reaches the user. Models that decline untrusted memory themselves
generate no prompts to approve.

Prior work on small open-weight models concluded that provenance cues visible to the model do not
hold~\cite{ppmf,sleeperchannels}. Our Llama result agrees. The hosted models show that this is a
property of the model rather than of labels: three of them treated a visible label as binding for
requests that act. A deployment cannot know in advance which kind of model it has, and a model
update can change it, so the case for enforcement stands. But for the three label-respecting models
enforcement left utility unchanged relative to labels alone and moved attack success by at most
two points.

\begin{table}[t]\centering\small
\TableFatigue
\caption{Enforced condition: attack success (\%) as the simulated user becomes more willing to
approve a confirmation prompt for a planted value.}
\label{tab:fatigue}
\end{table}

\subsection{What the defense costs when nobody is attacking}
\label{sec:cost}

\begin{table}[t]\centering\small
\resizebox{\textwidth}{!}{\TableCost}
\caption{Cost with no planted items. Share of requests with a known right answer. ``Detail lost'':
no retrievable entry held the right value. ``Asked anyway'': one did, and the model asked the
user instead. Mean $\pm$ standard deviation over five seeds.}
\label{tab:cost}
\end{table}

Table~\ref{tab:cost} gives the cost. With no attacker, the share of requests completed correctly
without involving the user fell by between 15 points (Haiku) and 43 points (Gemini) when the
enforced defense was switched on. Three mechanisms account for it.

\emph{The model asks.} Once Sonnet, Gemini and Haiku can see that a legitimate supplier detail is
labelled untrusted, they decline to act on it and ask the user: on 22\%, 30\% and 21\% of
requests respectively. For these models the runtime's confirmation prompt never fired, because
the model never proposed an untrusted value for a request that acts.

\emph{The runtime asks.} Llama never asked on its own, so the gate did the work, at about 10
confirmation prompts per 100 requests; Terra produced 7 and Luna under 2. A confirmation is
cheaper for the user than an open question, since it shows the value and its source, and an
approved value becomes trusted and is not asked about again.

\emph{Details expire.} Untrusted entries expire after 20 sessions, and a legitimate fact may be
needed long after it arrived. Sonnet, Gemini and Llama lost no details with no defense and 13\%,
13\% and 8\% with enforcement.

Counting a confirmed request as correct, Llama completes 92\% of requests under enforcement and
Terra 83\%, against 100\% and 93\% undefended. For the label-respecting models there is nothing
comparable to add back, because in our harness a question to the user is not followed by an
answer that the agent stores (Section~\ref{sec:limits}).

\subsection{Lifecycle: merging, expiry and repair}
\label{sec:lifecycle}

\emph{Merging across trust levels damages the user's own instructions.} In the Labels condition a
merge of the user's stated account with a planted note yields one entry labelled untrusted and
due to expire. Under attack, the share of needed details lost was higher with Labels than with
Enforced for every model, by 1 to 12 points (for example 23\% against 11\% for Sonnet, and 18\%
against 7\% for Llama). It is also a likely source of the residual override success in the Labels
column, since with the trusted entry gone there is nothing for the model to prefer; we did not
trace each case. Keeping trust
levels apart during consolidation removed the effect.

\emph{Expiry is a blunt instrument.} It bounds how long a planted item can wait, and it removed
legitimate outside facts at the rates given above. A confirmation that promotes a value to
trusted is what keeps a useful outside fact alive; Llama, which reaches the confirmation prompt
most often, lost fewer details under Enforced (8\%) than under Labels (13\%).

\emph{Repair depends on lineage.} At the end of each run we retract every document that carried a
planted item. With no defense, almost all the planted items that had been saved remained (the
rest had already been overwritten by a later merge), because the agent's notes about them
recorded no source to retract by. With labels, none remained for any
model.

\section{Limitations}
\label{sec:limits}

\begin{itemize}
\item \textbf{The attacks do not adapt.} The planted texts are fixed, simply worded templates.
Defenses that look strong against static attacks have repeatedly been broken by adaptive
ones~\cite{attackermovessecond}. Our undefended and Labels figures are therefore lower bounds on
attack success. The Enforced rules do not depend on wording, but the confirmation prompt depends
on a person, and we have not tested it with people.
\item \textbf{Synthetic workload, written by us.} Requests are short, topics are of three kinds,
and each document statement carries a topic handle that stands in for retrieval. Real memory
systems retrieve by similarity, which adds attack surface we do not model. Rates of asking,
prompting and expiry will move with the task mix, the expiry length and the share of facts that
arrive from outside.
\item \textbf{Asking is not followed by learning.} When the model asks the user for a detail, the
harness records the question and stores no answer, so the same detail can be asked about again.
An approved confirmation, by contrast, is remembered. This overstates the long-run burden for
models that ask relative to models that trigger confirmations, and should be corrected before
the comparison between the two is relied on.
\item \textbf{One run configuration per model.} Five seeds, one prompt wording, one expiry
length. We report standard deviations and no significance tests. Replies are not guaranteed to be
deterministic. Four models were reached through an aggregator and two through their vendor's API.
\item \textbf{Our implementation only.} We did not run the released implementations of the
defenses we cite. The Enforced condition is our own reading of those mechanisms.
\item \textbf{Empty replies.} Twenty-one replies from the two Claude models were empty and are
scored as saving nothing; the cause was not recorded.
\item \textbf{Not a model ranking.} A single synthetic workload cannot establish which model is
safer in general, and model behaviour changes between versions.
\end{itemize}

\section{Conclusion}

On this workload, provenance labels shown to the model were a strong defense for models that
respect them and a weak one for a model that does not, and enforcement by the runtime removed
that dependence. The enforced defense blocked every simple planted item we tried, at the price
of involving the user in 15 to 43 points more requests when nobody was attacking, through model
questions, runtime confirmations and expired facts. Which of those a deployment pays depends on
the model. The measurements that would most change this picture are adaptive attacks, a realistic
benchmark such as MemoryDojo, and a harness in which the user's answers are learned.

\section*{Disclosure}

An AI assistant (Claude, from Anthropic) wrote the harness and analysis code and drafted this
manuscript under the author's direction, and the author ran the experiments. Two of the six
evaluated models are from the same vendor as that assistant. All figures are computed by a
script from the saved transcripts, which are released so that they can be recounted.

\section*{Availability}

Code, workload generator, per-model reports and full transcripts:
\url{https://github.com/Radmo3/memgate} (private at the time of this draft).

\begin{thebibliography}{99}\small
% TODO before submission: complete author lists from the arXiv records. Entries without
% authors below are ones where the author list was not verified while drafting.
\bibitem{tmanm} Y.~Louck. Securing LLM-Agent Long-Term Memory Against Poisoning: Non-Malleable,
Origin-Bound Authority with Machine-Checked Guarantees. arXiv:2606.24322, 2026.
\bibitem{permit} PerMIT. arXiv:2610.28987, 2026.
\bibitem{ppmf} J.~Xu, Y.~Xiao, W.~Shao, H.~Liu, X.~Li. Memory Provenance Laundering in LLM Agents:
A Non-Amplification Firewall for Persistent Memory. arXiv:2607.29167, 2026.
\bibitem{sleeperchannels} N.~Maloyan, D.~Namiot. Sleeper Channels and Provenance Gates: Persistent
Prompt Injection in Always-on Autonomous AI Agents. arXiv:2605.13471, 2026.
\bibitem{priceofsafety} The Price of Safety: Benign-Case Utility and Token Overhead of
Memory-Poisoning Defenses in LLM Agents. arXiv:2609.22818, 2026.
\bibitem{memgauge} C.~Zang, Z.~Cao, X.~Meng, J.~Wang, W.~Chen, X.~Gao, L.~Wang, Z.~Li, S.~Guo.
Understanding Stage-Wise Utility-Risk Trade-offs in LLM Agent Memory. arXiv:2608.30177, 2026.
\bibitem{utilityunderattack} A.~Karunanidhi. Utility Under Attack: Agent Memory Poisoning and the
Limits of Content Screening and Provenance Ranking. arXiv:2608.21230, 2026.
\bibitem{trustgap} J.~Hu, S.~Ramachandran. The Memory Trust Gap: Capability-Dependent Failures in
Persistent-Memory Agents. arXiv:2609.01852, 2026.
\bibitem{rememberverify} B.~Li, J.~Yao, Z.~Zheng. Remember, Verify, or Ask? Cross-Family
Evaluation of Memory Commitment in LLM Agents. arXiv:2608.19564, 2026.
\bibitem{memorydojo} MemoryDojo. arXiv:2611.09721, 2026.
\bibitem{attackermovessecond} The Attacker Moves Second: Stronger Adaptive Attacks Bypass Defenses
Against LLM Jailbreaks and Prompt Injections. arXiv:2510.09023, 2025.
\end{thebibliography}

\end{document}
